\documentclass[11pt,a4paper]{article}
\usepackage[margin=2.5cm]{geometry}
\usepackage{amsmath,amsthm,amssymb}
\usepackage[hidelinks]{hyperref}
\usepackage{cleveref}
\usepackage[backend=biber,style=numeric,urldate=long]{biblatex}
\addbibresource{references.bib}

\theoremstyle{definition}
\newtheorem{definition}{Definition}[section]
\newtheorem{example}[definition]{Example}
\theoremstyle{plain}
\newtheorem{theorem}[definition]{Theorem}
\newtheorem{lemma}[definition]{Lemma}
\theoremstyle{remark}
\newtheorem{remark}[definition]{Remark}

\title{Walks, paths, and connectivity from the adjacency matrix: a short account}
\author{}
\date{\today}

\begin{document}
\maketitle

\begin{abstract}
For a finite simple undirected graph with adjacency matrix $A$, the entry
$(A^k)_{ij}$ is the number of walks of length $k$ from vertex $n_i$ to vertex $n_j$
(walks, not paths). A path is a walk with no repeated vertex. A graph on $N$ vertices is
connected if and only if every off-diagonal entry of $A+A^2+\dots+A^{N-1}$ is
positive. Each statement is proved and the test is run on two graphs.
\end{abstract}

\section{Introduction}

Let $G$ be a finite simple undirected graph with $N$ vertices $n_1,\dots,n_N$
and adjacency matrix $A$ (defined in \cref{sec:setting}). The lecture board
asks two questions about it, and the assignment adds a third
\cite{whiteboard1,whiteboard2,userstatements}.
\begin{enumerate}
\item[(Q1)] What does the entry $(A^k)_{ij}$ represent? The mark above $A$ on
  the board is read as an exponent $k$ \cite{userstatements}.
  \emph{Answer:} it is the number of walks of length $k$ from $n_i$ to $n_j$
  (\cref{thm:counting}); it counts walks, not paths.
\item[(Q2)] How can one tell from $A$ whether $G$ is connected?
  \emph{Answer:} $G$ is connected if and only if every entry off the
  diagonal of $B=A+A^2+\dots+A^{N-1}$ is positive (\cref{thm:sum}).
\item[(Q3)] What is a path between $n_i$ and $n_j$? \emph{Answer:} a walk
  from $n_i$ to $n_j$ in which no vertex occurs twice
  (\cref{def:pathbetween}).
\end{enumerate}
The answers come in the order Q1, Q3, Q2, because the test for Q2 needs the path. A full-text search of the course textbook
\cite{petrosyan2015graphs} finds neither the walk-counting theorem nor a matrix test for
connectivity, so both are taken from other sources, and every proof is written out in
full; ``derived'' marks this paper's own proof of a result stated in a cited source.
Textbook page numbers are printed ones, one lower than the PDF's.

Plain-words paragraphs precede the formal definitions, theorems, lemmas, and the worked example. They tell one story. A pond has \emph{islands}, and some pairs of islands have a \emph{bridge} between them (vertices and edges); two islands joined by a bridge are \emph{next door}. A frog makes \emph{trips} of \emph{hops}, each hop crossing one bridge, and it may hop back or land twice on an island (a walk); a \emph{fresh-island trip} lands on no island twice (a path), and the pond is \emph{in one piece} when the frog can get from any island to any other (connected). A square \emph{table} with a one in each cell whose two islands are next door, and a zero elsewhere, is the adjacency matrix.

\section{Setting and definitions}\label{sec:setting}

\paragraph{Setting.} $G=(V,E)$ is a graph in the sense of the textbook:
a set $V=\{n_1,\dots,n_N\}$ of vertices, $N\ge 1$, and a set $E$ of edges
\cite[Definition~1.1.1, p.~9]{petrosyan2015graphs}. It is finite, simple
(no edge from a vertex to itself, at most one edge between two vertices) and
undirected (the edge joining $u$ and $v$, written $uv$, is the same as $vu$)
\cite{userstatements}. The vertices $u$ and $v$ are \emph{adjacent} if
$uv\in E$ \cite[Definition~1.1.4, p.~11]{petrosyan2015graphs}. The indices
$i,j,r$ run over $1,\dots,N$; $k\ge 0$ is an integer.

\plainwords{} A pond has islands, and some pairs of islands have a bridge between them. Draw a square table with one row and one column for each island. Put a one in a cell when the two islands are next door, and a zero in every other cell.

\begin{definition}[Adjacency matrix]\label{def:matrix}
The \emph{adjacency matrix} of $G$ is the $N\times N$ matrix $A$ with
$A_{ij}=1$ if $n_i$ and $n_j$ are adjacent and $A_{ij}=0$ otherwise
\cite[p.~11]{petrosyan2015graphs}; the board states the same rule in other
letters \cite{whiteboard1,whiteboard2}.
\end{definition}
Its entries are $0$ or $1$. Matrix powers are the usual ones: $(XY)_{ij}=\sum_{r=1}^{N}X_{ir}Y_{rj}$,
$A^0=I$ (the identity matrix) and $A^{k+1}=A^kA$. Every entry of
$A^k$ is a non-negative integer, since sums and products of non-negative
integers are non-negative integers (\emph{used in} \cref{thm:sum}).

\plainwords{} A frog sits on an island and hops over a bridge to a next-door island. It may hop back, or land on the same island again. The islands it visits, in order, form a trip. A frog that never hops makes a trip with no hops.

\begin{definition}[Walk]\label{def:walk}
A \emph{walk} from $u_0$ to $u_k$ is a sequence of vertices
$u_0u_1\cdots u_k$ in which each vertex is adjacent to the one before it; its
\emph{length} is $k$, the number of edges used, each counted as often as it
is used. Vertices and edges may repeat. This is the textbook's Definition~1.2.10
\cite[p.~21]{petrosyan2015graphs}, which lists edges as well; in a simple graph
the vertices fix the edges. This paper also allows $k=0$: a single vertex is a walk of length $0$ from itself to
itself. This is the paper's own convention; the textbook's usage assumes it
\cite[pp.~23, 49]{petrosyan2015graphs}.
\end{definition}

\plainwords{} A fresh-island trip lands on no island twice. So it never crosses the same bridge twice either. Every fresh-island trip is a trip, but a trip may land on an island twice.

\begin{definition}[Path]\label{def:path}
A \emph{path} is a walk whose vertices are pairwise distinct. This is the
course convention \cite{userstatements}. It is the textbook's \emph{simple
path} \cite[Definition~1.2.13, p.~22]{petrosyan2015graphs}; the textbook's
own word ``path'' means a walk with no repeated edge, and this paper never
uses the word in that sense.
\end{definition}

\plainwords{} A pond is in one piece when the frog can get from any island to any other. The way may cross many bridges. It is enough to ask for a fresh-island trip between each two islands.

\begin{definition}[Connected]\label{def:connected}
$G$ is \emph{connected} if for any two vertices $u,v$ there is a path from
$u$ to $v$. The board states this
\cite{whiteboard1,whiteboard2}; the textbook's Definition~2.1.1
\cite[p.~48]{petrosyan2015graphs} is the same with the textbook's own ``path''. A vertex is joined to itself by the path
of length $0$, so only pairs of different vertices matter.
\end{definition}

\paragraph{Board graph.} The board's matrix has rows $0\,1\,0$, $1\,0\,1$,
$0\,1\,0$ \cite{whiteboard1,whiteboard2}:
\[
A=\begin{pmatrix}0&1&0\\1&0&1\\0&1&0\end{pmatrix}.
\]
It encodes the graph with edges $n_1n_2$ and $n_2n_3$, which is the example
graph throughout. The board's drawing shows a third line that the matrix
does not encode; the matrix defines the example.

\section{Answer to Q1: the entries of $A^k$ count walks}\label{sec:q1}

\plainwords{} Take the island table, and use as many copies of it as the trip has hops. Each cell then counts the trips with exactly that many hops from its row island to its column island. These trips may land on an island twice, so they are not all fresh-island trips.

\begin{theorem}[Walk counting]\label{thm:counting}
For every integer $k\ge 0$ and all $i,j$, the entry $(A^k)_{ij}$ is the number
of walks of length $k$ from $n_i$ to $n_j$.
\end{theorem}
The University of Galway notes state this for $k\ge 1$ with an induction
proof \cite[Theorem~1.1.3, p.~4]{galwaynotes}; the case $k=0$ is this paper's
convention. The proof is derived here.
\begin{proof}
By induction on $k$.

\emph{Base case $k=0$.} $A^0=I$. A walk of length $0$ from $n_i$ to $n_j$ is a
single vertex that is both start and end, so there is exactly one such walk
if $i=j$ and none if $i\ne j$. This is $I_{ij}$.

\emph{Inductive case.} Assume that for some $k\ge 0$ and all $i,r$ the entry
$(A^k)_{ir}$ is the number of walks of length $k$ from $n_i$ to $n_r$. Fix $i$
and $j$. Since $A^{k+1}=A^kA$,
\begin{equation}\label{eq:step}
(A^{k+1})_{ij}=\sum_{r=1}^{N}(A^k)_{ir}A_{rj}.
\end{equation}
A walk $u_0\cdots u_ku_{k+1}$ of length $k+1$ from $n_i$ to $n_j$ has
$u_k=n_r$ for exactly one $r$; sort the walks by this $r$. Fix $r$.
If $n_r$ is adjacent to $n_j$ (so $A_{rj}=1$), the walks of this kind
correspond one to one to the walks of length $k$ from $n_i$ to $n_r$: deleting
the final vertex $n_j$ gives such a walk, and appending $n_j$ to such a walk
gives a walk of the kind, because $n_r$ is adjacent to $n_j$. By the
hypothesis there are $(A^k)_{ir}=(A^k)_{ir}A_{rj}$ of them. If $n_r$ is not
adjacent to $n_j$ (so $A_{rj}=0$), no walk can end with the step from $n_r$ to
$n_j$, so there are $0=(A^k)_{ir}A_{rj}$ of them. Every walk is counted for
exactly one $r$, so the number of walks of length $k+1$ from $n_i$ to $n_j$ is
$\sum_r(A^k)_{ir}A_{rj}$, which is $(A^{k+1})_{ij}$ by \cref{eq:step}.
\end{proof}

\paragraph{Warning: walks, not paths.} The proof uses only that consecutive
vertices are adjacent, so repeated vertices are counted. For a path, appending
$n_j$ may repeat a vertex and the one to one correspondence fails. In the
board graph $(A^2)_{11}=1$ and the walk counted is $n_1n_2n_1$, which repeats
$n_1$; the only path from $n_1$ to $n_1$ is the single vertex $n_1$, of length
$0$ (a path of length $k\ge 1$ from a vertex to itself would repeat it). So
``$(A^2)_{11}$ counts the paths of length $2$ from $n_1$ to $n_1$'' is false,
although ``counts the walks'' is true. Also $(A^4)_{22}=4$ by script, so the entry is
not ``$1$ if $n_j$ can be reached from $n_i$''.

\section{Answer to Q3: the path between $n_i$ and $n_j$}\label{sec:q3}

\plainwords{} Choose two islands. A fresh-island trip between them starts on the first island and ends on the second. A frog that stays put makes such a trip from an island to itself.

\begin{definition}[Path between $n_i$ and $n_j$]\label{def:pathbetween}
A \emph{path between $n_i$ and $n_j$} (also a path \emph{from $n_i$ to $n_j$})
is a sequence $u_0,u_1,\dots,u_k$ with $k\ge 0$, $u_0=n_i$, $u_k=n_j$, in
which $u_{t-1}$ and $u_t$ are adjacent for $1\le t\le k$ and the vertices
$u_0,\dots,u_k$ are pairwise distinct. Its length is $k$.
\end{definition}
This is \cref{def:path} with the ends fixed, the textbook's
Definition~1.2.13 with ends $n_i$ and $n_j$
\cite[p.~22]{petrosyan2015graphs}. If $i\ne j$ such a path has at least one
edge; if $i=j$ the only path is $(n_i)$, since $u_0=u_k$ with $k\ge 1$ would
repeat a vertex. In the board graph $n_1n_2n_3$ is the only path from $n_1$ to $n_3$.

\section{Answer to Q2: a connectivity test from $A$}\label{sec:q2}

The test rests on two facts about walks and paths. Both are stated in
Grinberg \cite[Corollary~2.9.7, p.~39; Section~2.9.4, p.~40]{grinberg2025graph}
and the first is the textbook's Lemma~1.2.1 for $u\ne v$ without a length
bound \cite[p.~23]{petrosyan2015graphs}; the proofs are derived here.

\plainwords{} Suppose a frog gets from one island to another and lands on some island twice. Cut out all the hops between the two visits to that island. What is left is still a trip between the same two islands, with fewer hops. Cutting again and again ends in a fresh-island trip.

\begin{lemma}[Every walk contains a path]\label{lem:cut}
Let $W=u_0\cdots u_k$ be a walk from $u$ to $v$. Then there is a path from $u$
to $v$ of length at most $k$.
\end{lemma}
\begin{proof}
By strong induction on $k$: for walks of length $k$ we assume the statement
for all shorter walks. \emph{Base case $k=0$.} $W=(u)$ has one vertex, so it
is a path of length $0$. \emph{Inductive case $k\ge 1$.} If the vertices of $W$
are pairwise distinct, $W$ is a path of length $k$. Otherwise there are
positions $s<t$ with $u_s=u_t$. Let $W'=u_0\cdots u_su_{t+1}\cdots u_k$
(if $t=k$, $W'=u_0\cdots u_s$). It is a walk, because the only new
consecutive pair is $u_s,u_{t+1}$ and $u_s=u_t$ is adjacent to $u_{t+1}$. It
runs from $u$ to $v$ and has length $k-(t-s)<k$. By the hypothesis $W'$, and so
$W$, yields a path from $u$ to $v$ of length at most $k-(t-s)\le k$.
\end{proof}

\plainwords{} Pick two different islands and look at their cell in the tables for one hop, two hops, and so on. Go up to one fewer hops than there are islands. The frog can get from one island to the other exactly when one of these cells is not zero.

\begin{lemma}[Reachability from the powers]\label{lem:reach}
Let $i\ne j$. There is a path from $n_i$ to $n_j$ if and only if
$(A^k)_{ij}>0$ for some $k$ with $1\le k\le N-1$.
\end{lemma}
\begin{proof}
($\Rightarrow$) Let a path from $n_i$ to $n_j$ have length $k$. As $i\ne j$,
$k\ge 1$. A path of length $k$ has $k+1$ pairwise distinct vertices among
the $N$ vertices of the graph, so $k\le N-1$. A path is a walk, so by
\cref{thm:counting} $(A^k)_{ij}\ge 1$. ($\Leftarrow$) If $(A^k)_{ij}>0$ with
$1\le k\le N-1$, then by \cref{thm:counting} a walk of length $k$ joins $n_i$
to $n_j$, and by \cref{lem:cut} a path joins them.
\end{proof}

\plainwords{} Add the tables for one hop, two hops, and so on, up to one fewer hops than there are islands. Look only at cells for two different islands. The pond is in one piece exactly when none of these cells is zero. How large a cell is does not matter.

\begin{theorem}[Sum test]\label{thm:sum}
Let $N\ge 2$ (for $N=1$ the graph is connected) and $B=A+A^2+\dots+A^{N-1}$. Then $G$ is connected if and only if
$B_{ij}>0$ for all $i\ne j$.
\end{theorem}
This is Theorem~2 of Paramadevan and Sotheeswaran
\cite[p.~17]{paramadevan2021adjacency}; the proof here is derived.
\begin{proof}
Each $A^k$ has non-negative entries, so $B_{ij}=\sum_{k=1}^{N-1}(A^k)_{ij}$ is a
sum of non-negative terms, and it is positive exactly when at least one term
is positive. For $i\ne j$ this holds exactly when a path joins $n_i$ and
$n_j$, by \cref{lem:reach}. So $B_{ij}>0$ for all $i\ne j$ exactly when every
pair of different vertices is joined by a path, which is
\cref{def:connected}.
\end{proof}
The cutoff $N-1$ is the largest length a path can have
(proof of \cref{lem:reach}). The entries of $B$ are walk counts, not path counts; only
whether they are zero is used.

\section{Worked example and limits}\label{sec:example}

\plainwords{} We check the answer on two ponds with three islands. In the first, the islands lie in a row and the frog can get anywhere. In the second, two islands share a bridge and the third island stands alone.

All matrices and walk lists below were computed by the scripts in the
\texttt{scripts/} folder, which also holds their outputs.

\paragraph{Connected: the board graph} ($N=3$, edges $n_1n_2$, $n_2n_3$).
Here $B=A+A^2$ with
\[
A^2=\begin{pmatrix}1&0&1\\0&2&0\\1&0&1\end{pmatrix},\qquad
B=\begin{pmatrix}1&1&1\\1&2&1\\1&1&1\end{pmatrix}.
\]
By \cref{thm:counting}, $(A^2)_{11}=1$ counts the walk
$n_1n_2n_1$, $(A^2)_{13}=1$ counts $n_1n_2n_3$, $(A^2)_{22}=2$ counts
$n_2n_1n_2$ and $n_2n_3n_2$, and $(A^2)_{12}=0$ because no walk of length $2$
joins $n_1$ to $n_2$. The script also gives $(A^3)_{12}=2$, for the walks
$n_1n_2n_1n_2$ and $n_1n_2n_3n_2$, and $(A^4)_{22}=4$. All off-diagonal entries of $B$ are positive, so $G$ is
connected by \cref{thm:sum}: the pairs $n_1n_2$ and $n_2n_3$ get a positive
entry from $A$, and the pair $n_1n_3$ gets $B_{13}=1$ from $A^2$, which counts
the path $n_1n_2n_3$.

\paragraph{Disconnected: the edge $n_1n_2$ and the isolated vertex $n_3$}
($N=3$). Then
\[
A=\begin{pmatrix}0&1&0\\1&0&0\\0&0&0\end{pmatrix},\quad
A^2=\begin{pmatrix}1&0&0\\0&1&0\\0&0&0\end{pmatrix},\quad
B=\begin{pmatrix}1&1&0\\1&1&0\\0&0&0\end{pmatrix}.
\]
The entries $B_{13}=B_{23}=B_{31}=B_{32}=0$, so $G$ is not connected by
\cref{thm:sum}. This agrees with the graph: no walk joins $n_3$ to $n_1$ or $n_2$.

\paragraph{Limits.} The test answers the question using $A$ alone, but it has
limits. It reads only whether entries are zero: a positive entry says a walk of
exactly length $k$ exists, hence a path of length at most $k$, not a path of
length exactly $k$. The diagonal of $B$ is ignored, and a connected graph can
have zeros there: for $N=2$ and one edge, $B=A$ has a zero diagonal. The sum must
run to $A^{N-1}$: on the path graph with edges $n_1n_2$, $n_2n_3$, $n_3n_4$
($N=4$) the script gives $(A^1)_{14}=(A^2)_{14}=0$ and $(A^3)_{14}=1$, so a sum
stopped at $A^2$ would call this connected graph disconnected. The test is proved only for finite simple undirected graphs.


\printbibliography

\end{document}
